\documentclass{article}



\usepackage{arxiv}

\usepackage[utf8]{inputenc} % allow utf-8 input
\usepackage[T1]{fontenc}    % use 8-bit T1 fonts
\usepackage{hyperref}       % hyperlinks
\usepackage{url}            % simple URL typesetting
\usepackage{booktabs}       % professional-quality tables
\usepackage{amsfonts}       % blackboard math symbols
\usepackage{nicefrac}       % compact symbols for 1/2, etc.
\usepackage{microtype}      % microtypography
\usepackage{lipsum}		% Can be removed after putting your text content
\usepackage{graphicx}
\usepackage{natbib}
\usepackage{doi}



\title{Proposal for refinements to the metric system\\as used in astronomy}

\date{April 1th, 2022}	% Here you can change the date presented in the paper title
%\date{} 					% Or removing it

\author{ \href{https://orcid.org/0000-0001-9527-436X}{\includegraphics[scale=0.06]{orcid.pdf}\hspace{1mm}S.P.C.~Peters}\\
	Department of Computer Science\\
	Windesheim University OAS\\
	Zwolle, The Netherlands \\
	\texttt{s.p.c.peters@windesheim.nl} \\
	%% examples of more authors
	\And
	\href{https://orcid.org/0000-0002-6033-960X}{\includegraphics[scale=0.06]{orcid.pdf}\hspace{1mm}E.G.P.~Bos} \\
	Netherlands eScience Center\\
	Amsterdam, The Netherlands\\
	\texttt{p.bos@esciencecenter.nl} \\
	\And
	\href{https://orcid.org/0000-0002-6033-960X}{\includegraphics[scale=0.06]{orcid.pdf}\hspace{1mm}Y.G.~Grange} \\
	ASTRON\\
	Netherlands Institute for Radio Astronomy\\
	Dwingeloo, The Netherlands\\
	\texttt{grange@astron.nl} \\
	\And
	\href{https://orcid.org/0000-0002-6033-960X}{\includegraphics[scale=0.06]{orcid.pdf}\hspace{1mm}E.~Busekool} \\
	SG De Waerdenborch\\
	Ommen, The Netherlands\\
	\texttt{?@?.com} \\
}

% Uncomment to remove the date
%\date{}

% Uncomment to override  the `A preprint' in the header
%\renewcommand{\headeright}{Technical Report}
%\renewcommand{\undertitle}{Technical Report}
\renewcommand{\shorttitle}{Let's make things better}

%%% Add PDF metadata to help others organize their library
%%% Once the PDF is generated, you can check the metadata with
%%% $ pdfinfo template.pdf
\hypersetup{
pdftitle={A template for the arxiv style},
pdfsubject={q-bio.NC, q-bio.QM},
pdfauthor={David S.~Hippocampus, Elias D.~Striatum},
pdfkeywords={First keyword, Second keyword, More},
}

\begin{document}
\maketitle

\begin{abstract}
	We review the current status of the various metric systems which are available on the market. 
	In particular, we focus on their application to the field of astronomy.
	The various advantages and disadvantages inherent to each system are reviewed.
	A new unifying system of metrics is subsequently proposed.
	It is demonstrated that the new system outperforms the previously available systems by a significant margin. 
	In particular, we demonstrate that the new system is superior at captivating the minds of the public.
\end{abstract}


% keywords can be removed
\keywords{Metric system \and Proposal \and Giraffes}


\section{Introduction}
Astronomy has always been a field of extremes. On one end, astronomers measure properties of large objects that are very far away \textbf{blatant self-citations here} while on the other hand, units are used that represent relatively small ranges \textbf{other blatant self-citations to papers in cgs and angstrom go here}, or units (e.g. parsec, solar mass) that may seem practical at first, but the human intuition is unable to handle those. Several times, it has been suggested that the world of astronomy change their way to use units that are part of the International system of units, but this has never resulted in a wide adoption. 

Taking this into account, the thing the astronomical world needs most is a system of units that is of a higher magnitude than the SI, while still being numbers that people can parse. It is widely known that people are very bad at having intuition for non-tangible objects. It is therefore essential for a scientific system of units to actually be connected to the real-life experience of the average human being. 

\section{proposing the Giraffal measurement system}
\begin{table}
\begin{tabular}{lll}
\hline
Unit & 1 Giraffe \\
\hline
\hline
cm & 600 \\
\hline
m & 6 \\
\hline 
football field & 17.5 in the UK, 18 in the US, 22.8 in Canada, 22.5--31 in Australia and 22-24 in Ireland. \\
\hline
parsec & big number \\
\hline
\hline
\end{tabular}
\caption{•}
\end{table}

An added advantage of the giraffic system is the application of the speed of light. Since the speed of light is generally expressed in the SI system, we first need to convert the standard time and length from the giraffic system to SI. Since one giraffe is 6m, the period of a 1G pendulum can easily be calculated using the equation

\begin{equation}
 T = 2\pi \sqrt{\frac{L}{g}}
\end{equation}
This leads to a Giraffe time of ~5 seconds. On the other hand, $3\times10^{5}\,km$ can be expressed as $5\times10^{4}\,Gl$. Dividing this leads to a light speed of $1\times10^{4} GirL/GirT$. It is obvious that when a unit system is defined in such a way that constants of nature become 1, this system of units is by definition superior. Also, using the Giraf in ones publication makes the integration between particle physics (a field that clasically has been using a light speed of 1) and astronomy far easier. 

\section{Headings: second level}


\section{Headings: third level}


\section{Conclusion}

\section{Discussion}
TODO: rebuttal van the elephantic system





\subsection{Citations}
Citations use \verb+natbib+. The documentation may be found at
\begin{center}
	\url{http://mirrors.ctan.org/macros/latex/contrib/natbib/natnotes.pdf}
\end{center}

Here is an example usage of the two main commands (\verb+citet+ and \verb+citep+): Some people thought a thing \citep{kour2014real, hadash2018estimate} but other people thought something else \citep{kour2014fast}. Many people have speculated that if we knew exactly why \citet{kour2014fast} thought this\dots

\subsection{Figures}
\lipsum[10]
See Figure \ref{fig:fig1}. Here is how you add footnotes. \footnote{Sample of the first footnote.}
\lipsum[11]

\begin{figure}
	\centering
	\fbox{\rule[-.5cm]{4cm}{4cm} \rule[-.5cm]{4cm}{0cm}}
	\caption{Sample figure caption.}
	\label{fig:fig1}
\end{figure}

\subsection{Tables}
See awesome Table~\ref{tab:table}.

The documentation for \verb+booktabs+ (`Publication quality tables in LaTeX') is available from:
\begin{center}
	\url{https://www.ctan.org/pkg/booktabs}
\end{center}


\begin{table}
	\caption{Sample table title}
	\centering
	\begin{tabular}{lll}
		\toprule
		\multicolumn{2}{c}{Part}                   \\
		\cmidrule(r){1-2}
		Name     & Description     & Size ($\mu$m) \\
		\midrule
		Dendrite & Input terminal  & $\sim$100     \\
		Axon     & Output terminal & $\sim$10      \\
		Soma     & Cell body       & up to $10^6$  \\
		\bottomrule
	\end{tabular}
	\label{tab:table}
\end{table}

\subsection{Lists}
\begin{itemize}
	\item Lorem ipsum dolor sit amet
	\item consectetur adipiscing elit.
	\item Aliquam dignissim blandit est, in dictum tortor gravida eget. In ac rutrum magna.
\end{itemize}


\bibliographystyle{unsrtnat}
\bibliography{references}  %%% Uncomment this line and comment out the ``thebibliography'' section below to use the external .bib file (using bibtex) .


%%% Uncomment this section and comment out the \bibliography{references} line above to use inline references.
% \begin{thebibliography}{1}

% 	\bibitem{kour2014real}
% 	George Kour and Raid Saabne.
% 	\newblock Real-time segmentation of on-line handwritten arabic script.
% 	\newblock In {\em Frontiers in Handwriting Recognition (ICFHR), 2014 14th
% 			International Conference on}, pages 417--422. IEEE, 2014.

% 	\bibitem{kour2014fast}
% 	George Kour and Raid Saabne.
% 	\newblock Fast classification of handwritten on-line arabic characters.
% 	\newblock In {\em Soft Computing and Pattern Recognition (SoCPaR), 2014 6th
% 			International Conference of}, pages 312--318. IEEE, 2014.

% 	\bibitem{hadash2018estimate}
% 	Guy Hadash, Einat Kermany, Boaz Carmeli, Ofer Lavi, George Kour, and Alon
% 	Jacovi.
% 	\newblock Estimate and replace: A novel approach to integrating deep neural
% 	networks with existing applications.
% 	\newblock {\em arXiv preprint arXiv:1804.09028}, 2018.

% \end{thebibliography}


\end{document}
